\documentclass[UTF8]{ctexart}
\usepackage{amsmath, amssymb, amsthm}
\usepackage{geometry}
\usepackage{xcolor}
\usepackage{tcolorbox}
\usepackage{enumitem}
\usepackage{fancyhdr}
\usepackage{booktabs}
\usepackage{tikz} % 用于绘图
\tcbuselibrary{skins, breakable}

% 页面设置
\geometry{a4paper, margin=1in}
\pagestyle{fancy}
\fancyhead[L]{专升本高数：微分中值定理考点全覆盖}
\fancyhead[R]{劳资必是本科生}

% ==========================================
% 环境定义区
% ==========================================

% 1. 基础红色字体环境（符合讲义要求）
\newenvironment{redenv}{\color{red}}{}

% 2. 概念解析框（理论部分用，也可用于习题概念）
\newtcolorbox{conceptbox}[1][概念解析]{
    colback=red!5, 
    colframe=red!50!black, 
    fonttitle=\bfseries, 
    title=#1, 
    breakable
}

% 3. 真题提示框（红色系，醒目）
\newtcolorbox{examnote}[1]{
    colback=red!5!white,
    colframe=red!75!black,
    fonttitle=\bfseries,
    title=【真题提示：#1】,
    breakable
}

% 4. 详细解析专用环境（习题部分用）
% 学生版：定义为空或只保留标题占位，但为了排版美观，我们在正文中直接用 \vspace 替换内容
\newenvironment{detailedanalysis}{%
    \color{red}
    \begin{enumerate}[label=\textbf{第\arabic*步：}, leftmargin=*, align=left, itemsep=1.5ex]
}{%
    \end{enumerate}
}

% 5. 公式推导环境（习题部分用）
\newenvironment{formuladerivation}{%
    \begin{enumerate}[label={\textbf{公式\arabic*：}}, leftmargin=*, align=left]
}{%
    \end{enumerate}
}

% 6. 分隔线命令
\newcommand{\problemrule}{\noindent\rule{\textwidth}{1.5pt}\vspace{-0.8em}}
\newcommand{\analysisrule}{\noindent\rule{\textwidth}{0.8pt}\vspace{-0.8em}}

\begin{document}

\title{\textbf{微分中值定理：方法总结与归纳（学生版）}}
\author{fifine老师}
\date{}
\maketitle

% ##################################################################
% 第一部分：理论体系与逻辑架构 (讲义内容)
% ##################################################################

\section{引言：局部性态向整体特征的逻辑跨越}
微分中值定理是连接函数“局部性质”与“整体形态”的桥梁，也是解决专升本高数证明题的核心钥匙。

\textbf{通俗理解：}
如果我们只知道函数在两端 $a$ 和 $b$ 的取值（这是整体信息），通常很难直接看出函数中间的具体变化。但是，中值定理告诉我们：\textbf{整体的平均变化率，一定等于中间某一点的瞬时变化率。}

\textbf{物理直观：}
如果你开车 2 小时走了 120 公里（整体平均速度 $60\text{km/h}$），那么在途中一定有一个瞬间，你的车速表显示的正是 $60\text{km/h}$（瞬时速度）。这就是“以瞬时（导数）刻画整体（差商）”的核心思想。

\textbf{解题意义：}
它允许我们将难处理的函数差值 $f(b)-f(a)$，转化为某一点的导数 $f'(\xi)$ 来研究。这样，我们就可以利用导数的工具（正负号判断单调性等）来解决关于原函数的证明问题。

\subsection*{核心定理}

\begin{figure}[h!]
    \centering
    \begin{minipage}{0.48\textwidth}
        \centering
        \begin{tikzpicture}[scale=0.8, domain=0.5:3.5]
            \draw[->] (-0.2,0) -- (4,0) node[right] {$x$};
            \draw[->] (0,-0.2) -- (0,3.5) node[above] {$y$};
            \draw[dashed] (1,0) node[below]{$a$} -- (1,2);
            \draw[dashed] (3,0) node[below]{$b$} -- (3,2);
            \draw[thick, blue] plot (\x, {-1.0*(\x-2)*(\x-2) + 3});
            \draw[thick, red] (1.2, 3) -- (2.8, 3);
            \fill (2,3) circle (2pt) node[above] {$\xi$};
            \draw[dashed] (2,0) node[below]{$\xi$} -- (2,3);
            \node at (2, -1) {\textbf{罗尔定理} $f'(\xi)=0$};
        \end{tikzpicture}
        \begin{tcolorbox}[colback=red!5, colframe=red!75!black, title=罗尔定理定义]
        若 $f \in C[a,b]$，$(a,b)$ 可导，且 \textbf{$f(a)=f(b)$}，\\
        则 $\exists \xi \in (a,b)$，使得 \textbf{$f'(\xi)=0$}。
        \end{tcolorbox}
    \end{minipage}
    \hfill
    \begin{minipage}{0.48\textwidth}
        \centering
        \begin{tikzpicture}[scale=0.8, domain=0.5:3.5]
            \draw[->] (-0.2,0) -- (4,0) node[right] {$x$};
            \draw[->] (0,-0.2) -- (0,4) node[above] {$y$};
            \draw[dashed] (1,0) node[below]{$a$} -- (1,1.25);
            \draw[dashed] (3,0) node[below]{$b$} -- (3,3.25);
            \draw[gray, dashed] (1,1.25) -- (3,3.25); % 割线
            \draw[thick, blue] plot (\x, {0.25*\x*\x + 1}); % 曲线
            \draw[thick, red] (1,1) -- (3,3); % 平行切线
            \fill (2,2) circle (2pt) node[below right] {$\xi$};
            \draw[dashed] (2,0) node[below]{$\xi$} -- (2,2);
            \node at (2, -1) {\textbf{拉格朗日定理} $f'(\xi)=\frac{f(b)-f(a)}{b-a}$};
        \end{tikzpicture}
        \begin{tcolorbox}[colback=blue!5, colframe=blue!75!black, title=拉格朗日定理由来]
        若 $f \in C[a,b]$，$(a,b)$ 可导，\\
        则 $\exists \xi \in (a,b)$，使得 \textbf{$f(b)-f(a) = f'(\xi)(b-a)$}。
        \end{tcolorbox}
    \end{minipage}
\end{figure}

\subsection{中值定理-归纳图}
\begin{center}
\begin{tikzpicture}[node distance=2.5cm, auto]
    \node (R) [draw, rectangle, rounded corners, minimum height=1cm, align=center] {罗尔定理 \\ $f'(\xi)=0$};
    \node (L) [draw, rectangle, rounded corners, below left of=R, node distance=5cm, minimum height=1cm, align=center] {拉格朗日定理 \\ $f'(\xi)=\frac{f(b)-f(a)}{b-a}$};
    \node (C) [draw, rectangle, rounded corners, below right of=R, node distance=5cm, minimum height=1cm, align=center] {柯西定理 \\ (拉格朗日推广)};
    
    \draw[->, thick] (L) -- (R) node[midway, above, sloped] {特例：$f(a)=f(b)$};
    \draw[->, thick] (R) -- (L) node[midway, below, sloped] {构造辅助函数变形};
    \draw[->, dashed] (L) -- (C) node[midway, below] {推广};
\end{tikzpicture}
\end{center}

\subsection*{定理选用策略：题目特征判别}
在解题过程中，根据题目给出的条件特征（如区间性质、端点值），快速锁定使用的定理：

\begin{center}
\begin{tabular}{p{0.45\textwidth}|p{0.45\textwidth}}
\toprule
\textbf{出现以下特征 $\to$ 优先用罗尔定理} & \textbf{出现以下特征 $\to$ 优先用拉格朗日} \\
\midrule
1. \textbf{结论形式}：$f'(\xi)=0$ 或含 $f(\xi), f'(\xi)$ 的等式（如 $f'(\xi)+f(\xi)=0$） & 1. \textbf{结论形式}：含 $f(b)-f(a)$ 或此形式的变体 \\
2. \textbf{端点条件}：题目明确给出（或隐含）$f(a)=f(b)$ & 2. \textbf{不等式证明}：需证明 $f(b)>f(a)$ 或与导数大小相关的不等式 \\
3. \textbf{多次求导}：高阶导数问题（如 $f^{(n)}(\xi)=0$）通常多次用罗尔 & 3. \textbf{极限计算}：出现 $f(x)-f(y)$ 的极限型问题 \\
\bottomrule
\end{tabular}
\end{center}

\begin{tcolorbox}[colback=yellow!10, colframe=orange!80!black, title=【冲刺口诀】]
\textbf{闭连续，开可导，端点相等罗尔找；} \\
\textbf{端点不等差商造，拉格朗日能通考。}
\end{tcolorbox}

\section{罗尔定理与辅助函数构造}

\begin{conceptbox}[定理 1：罗尔定理]
若函数 $f(x)$ 满足以下条件：
1. 在闭区间 $[a, b]$ 上连续；
2. 在开区间 $(a, b)$ 内可导；
3. \textbf{端点值相等} $f(a)=f(b)$；\\
则必存在 $\xi \in (a, b)$，使得 $f'(\xi)=0$。
\end{conceptbox}

\subsection{构造辅助函数 $F(x)$ 的三大方法}
在证明题中，能否正确构造 $F(x)$ 是解题的逻辑核心。：

\begin{enumerate}
    \item \textbf{方法一：一阶线性微分方程逆推（公式法）}
    针对结论形式为 $f'(\xi) + g(\xi)f(\xi) = 0$。根据一阶线性微分方程的积分因子求解思路，构造：
    $$F(x) = f(x) e^{\int g(x) dx}$$
    
    \begin{examnote}{2017年真题}
    \textbf{题目}：证明 $f'(\xi) - \frac{3}{1-\xi}f(\xi) = 0$。\\
    \textbf{逻辑构造}：\vspace{1.5cm} % 学生填写区域
    \end{examnote}

    \item \textbf{方法二：莱布尼茨乘积法则（四则运算型）}
    针对结论形式为 $u' v + u v' = 0$。利用积的求导公式逆推，构造 $F(x) = u(x)v(x)$。补充：加法用乘，减法用除!
    
    \begin{examnote}{2018年真题}
    \textbf{题目}：证明 $f(\xi) + \xi f'(\xi) = 0$。\\
    \textbf{逻辑构造}：\vspace{1.5cm} % 学生填写区域
    \end{examnote}

    \item \textbf{方法三：直接积分原函数型}
    针对结论不含 $f(x)$ 且导数项独立的等式。
    
    \begin{examnote}{示例题目}
    \textbf{题目}：证明 $f'(\xi) = 1$。\\
    \textbf{逻辑构造}：\vspace{1.5cm} % 学生填写区域
    \end{examnote}
\end{enumerate}

\section{拉格朗日中值定理}

\begin{conceptbox}[定理 2：拉格朗日中值定理]
若 $f(x)$ 在闭区间 $[a, b]$ 上连续，且在开区间 $(a, b)$ 内可导，则至少存在一点 $\xi \in (a, b)$，使：
$$f(b) - f(a) = f'(\xi)(b - a)$$
\end{conceptbox}

\subsection{逻辑应用：不等式与极限的解析}

\begin{examnote}{2018年真题}
\textbf{真题考点}：证明 $x < \tan x < \frac{x}{\cos^2 x}$ （其中 $x \in (0, \frac{\pi}{2})$）。\\
\textbf{逻辑推演}：\vspace{2cm} % 学生填写区域
\end{examnote}

\begin{examnote}{2016年真题}
\textbf{真题考点}：证明 $|\arcsin x_1 - \arcsin x_2| \ge |x_1 - x_2|$。\\
\textbf{逻辑推演}：\vspace{2cm} % 学生填写区域
\end{examnote}

\section{高阶技巧：罗尔定理与零点定理的嵌套应用}
当端点值不直接相等（如 $F(a) \neq F(b)$）时，需运用\textbf{连续函数的零点定理}来寻找中间点。

\begin{examnote}{示例题目}
\textbf{逻辑模型}：已知 $f(0)=f(1)=0, f(\frac{1}{2})=1$，证明 $f'(\xi)=1$。\\
\textbf{证明逻辑}：\vspace{3cm} % 学生填写区域
\end{examnote}

\section{中值定理求解极限技巧}
\begin{examnote}{示例题目}
\textbf{高阶计算}：求极限 $\lim_{x \to 0} \frac{e^x - e^{\sin x}}{x^3}$。\\
\textbf{解析逻辑}：\vspace{3cm} % 学生填写区域
\end{examnote}

% ##################################################################
% 第二部分：典型习题深度解析 (习题内容，按要求补充)
% ##################################################################
\newpage
\section{典型习题深度解析}

% ========== 第1题 ==========
\problemrule
\subsection*{【题目1】罗尔定理与辅助函数构造实战}
已知函数 $f(x)$ 在区间 $[a, b]$ 上连续，在 $(a, b)$ 内可导。
请证明：若需证 $f'(\xi) - \frac{3}{1-\xi}f(\xi) = 0$ （其中 $\xi \in (a, b)$），应如何构造辅助函数？结合2017年真题思路，写出完整的构造逻辑。

\begin{redenv}
\subsubsection*{逐步解析}
\vspace{3cm} % 学生作答区域
\end{redenv}

% ========== 第2题 ==========
\vspace{2em}
\problemrule
\subsection*{【题目2】拉格朗日中值定理证明不等式实战}
证明：当 $x \in (0, \frac{\pi}{2})$ 时，有 $x < \tan x < \frac{x}{\cos^2 x}$。（源自2018年真题考点）

\begin{redenv}
\subsubsection*{逐步解析}
\vspace{8cm} % 学生作答区域
\end{redenv}

% ========== 第3题 ==========
\vspace{2em}
\problemrule
\subsection*{【题目3】中值定理与零点定理的综合应用}
已知函数 $f(x)$ 在 $[0, 1]$ 上二阶可导，且 $f(0)=f(1)=0, f(\frac{1}{2})=1$。
证明：存在 $\xi \in (0, 1)$，使得 $f'(\xi) = 1$。

\begin{redenv}
\subsubsection*{逐步解析}
\vspace{8cm} % 学生作答区域
\end{redenv}

% ========== 第4题 ==========
\vspace{2em}
\problemrule
\subsection*{【题目4】利用中值定理求极限}
求极限：$\lim_{x \to 0} \frac{e^x - e^{\sin x}}{x^3}$。

\begin{redenv}
\subsubsection*{逐步解析}
\vspace{8cm} % 学生作答区域
\end{redenv}

% ========== 第5题 ==========
\vspace{2em}
\problemrule
\subsection*{【题目5】拉格朗日中值定理证明双向不等式}
证明：当 $0 < a < b$ 时，有 $\frac{b-a}{b} < \ln \frac{b}{a} < \frac{b-a}{a}$。

\begin{redenv}
\subsubsection*{逐步解析}
\vspace{8cm} % 学生作答区域
\end{redenv}

% ========== 第6题 ==========
\vspace{2em}
\problemrule
\subsection*{【题目6】零点定理与单调性判别根的个数}
证明：方程 $x^3 + x - 1 = 0$ 在区间 $(0, 1)$ 内有且仅有一个实根。

\begin{redenv}
\analysisrule
\subsubsection*{逐步解析}
\vspace{6cm} % 学生作答区域
\end{redenv}

% ========== 第7题 ==========
\vspace{2em}
\problemrule
\subsection*{【题目7】辅助函数构造（乘积型）}
已知函数 $f(x)$ 在 $[0, 1]$ 上连续，在 $(0, 1)$ 内可导，且 $f(0)=f(1)=0$。证明：存在 $\xi \in (0, 1)$，使得 $f'(\xi) + 2\xi f(\xi) = 0$。

\begin{redenv}
\analysisrule
\subsubsection*{逐步解析}
\vspace{8cm} % 学生作答区域
\end{redenv}

% ========== 第8题 ==========
\vspace{2em}
\problemrule
\subsection*{【题目8】利用导数性质比较大小}
比较 $e^\pi$ 与 $\pi^e$ 的大小。

\begin{redenv}
\analysisrule
\subsubsection*{逐步解析}
\vspace{6cm} % 学生作答区域
\end{redenv}

% ========== 第9题 ==========
\vspace{2em}
\problemrule
\subsection*{【题目9】利用拉格朗日中值定理证明不等式}
证明：当 $x>0$ 时，$\frac{x}{1+x} < \ln(1+x) < x$。

\begin{redenv}
\analysisrule
\subsubsection*{逐步解析}
\vspace{8cm} % 学生作答区域
\end{redenv}

% ========== 第10题 ==========
\vspace{2em}
\problemrule
\subsection*{【题目10】简单的拉格朗日中值定理应用}
已知函数 $f(x)$ 在 $[0, 1]$ 上连续，在 $(0, 1)$ 内可导，且 $f(0)=0, f(1)=1$。证明：在 $(0, 1)$ 内至少有一点 $\xi$ 使 $f'(\xi)=1$。

\begin{redenv}
\analysisrule
\subsubsection*{逐步解析}
\vspace{6cm} % 学生作答区域
\end{redenv}

% ========== 第11题 ==========
\vspace{2em}
\problemrule
\subsection*{【题目11】辅助函数构造（幂函数型）}
设函数 $f(x)$ 在 $[1, 2]$ 上连续，在 $(1, 2)$ 内可导，且 $f(1) = 4f(2)$。证明：存在一点 $\xi \in (1, 2)$，使得 $2f(\xi) + \xi f'(\xi) = 0$。

\begin{redenv}
\analysisrule
\subsubsection*{逐步解析}
\vspace{6cm} % 学生作答区域
\end{redenv}

% ========== 第12题 ==========
\vspace{2em}
\problemrule
\subsection*{【题目12】利用单调性证明不等式}
证明：当 $x > 0$ 时，$\ln(1+x) > x - \frac{1}{2}x^2$。

\begin{redenv}
\analysisrule
\subsubsection*{逐步解析}
\vspace{8cm} % 学生作答区域
\end{redenv}

% ========== 第13题 ==========
\vspace{2em}
\problemrule
\subsection*{【题目13】拉格朗日中值定理计算题}
求函数 $f(x) = x^2 + 2x$ 在区间 $[0, 1]$ 上满足拉格朗日中值定理条件的 $\xi$ 值。

\begin{redenv}
\analysisrule
\subsubsection*{逐步解析}
\vspace{8cm} % 学生作答区域
\end{redenv}

% ========== 第14题 ==========
\vspace{2em}
\problemrule
\subsection*{【题目14】利用单调性证明正弦不等式}
证明：当 $x > 0$ 时，$\sin x < x$。

\begin{redenv}
\analysisrule
\subsubsection*{逐步解析}
\vspace{8cm} % 学生作答区域
\end{redenv}

% ========== 第15题 ==========
\vspace{2em}
\problemrule
\subsection*{【题目15】反复应用罗尔定理}
已知 $f(x)$ 在 $[1, 3]$ 上二阶可导，且 $f(1)=f(2)=f(3)=0$。证明：存在 $\xi \in (1, 3)$，使得 $f''(\xi)=0$。

\begin{redenv}
\analysisrule
\subsubsection*{逐步解析}
\vspace{8cm} % 学生作答区域
\end{redenv}

% ========== 第16题 ==========
\vspace{2em}
\problemrule
\subsection*{【题目16】单调性与零点判定}
证明：方程 $x^3 - 3x + 1 = 0$ 在区间 $(-2, 2)$ 内恰有三个实根。

\begin{redenv}
\analysisrule
\subsubsection*{逐步解析}
\vspace{8cm} % 学生作答区域
\end{redenv}

\end{document}
